% Practice Test 23 — Minnesota Edition
% Questions: 30 (15 MC + 15 SA)
% Generated by scripts/generate_practice_tests.py

\begin{practiceQuestion}{1-3-q05}{mc}
\begin{questionText}
A car travels at a constant speed. It goes $150$ miles in $2.5$ hours. What is the constant of proportionality?
\end{questionText}
\choiceA{$30$ mph}
\choiceB{$60$ mph}
\choiceC{$75$ mph}
\choiceD{$150$ mph}
\correctAnswer{B}
\explanation{$k = \frac{150}{2.5} = 60$ miles per hour.}
\end{practiceQuestion}

\begin{practiceQuestion}{1-5-q30}{gsa}
\begin{questionText}
Two friends start walking at the same time. The graph shows their distances. How many miles apart are they after $6$ hours?

\begin{center}
\begin{tikzpicture}[scale=0.7]
  \draw[step=1,gray!20,thin] (0,0) grid (7,8);
  \draw[thick,->] (0,0) -- (7.5,0) node[right]{\small Hours};
  \draw[thick,->] (0,0) -- (0,8.5) node[above]{\small Miles};
  \foreach \x in {1,...,7} \draw (\x,0.07) -- (\x,-0.07) node[below]{\tiny \x};
  \foreach \y in {1,...,8} \draw (0.07,\y) -- (-0.07,\y) node[left]{\tiny \y};
  \draw[thick,funBlue] (0,0) -- (6,7.2);
  \fill[funBlue] (0,0) circle (2pt);
  \fill[funBlue] (5,6) circle (2.5pt) node[above left]{\small $(5, 6)$};
  \node[above,funBlue] at (5.5,6.6) {\small Ava};
  \draw[thick,funRed] (0,0) -- (6,4.8);
  \fill[funRed] (0,0) circle (2pt);
  \fill[funRed] (5,4) circle (2.5pt) node[below right]{\small $(5, 4)$};
  \node[below right,funRed] at (5.5,4.4) {\small Ben};
\end{tikzpicture}
\end{center}
\end{questionText}
\correctAnswer{$2.4$ miles}
\explanation{Ava: $k = \frac{6}{5} = 1.2$ mph; at $6$ hr: $1.2 \times 6 = 7.2$ mi. Ben: $k = \frac{4}{5} = 0.8$ mph; at $6$ hr: $0.8 \times 6 = 4.8$ mi. Difference: $7.2 - 4.8 = 2.4$ miles.}
\end{practiceQuestion}

\begin{practiceQuestion}{1-6-q21}{sa}
\begin{questionText}
A scale drawing uses $1$ cm $= 4.5$ m. If a building is $27$ m tall, how tall is it in the drawing?
\end{questionText}
\correctAnswer{$6$ cm}
\explanation{$\frac{1}{4.5} = \frac{x}{27}$. Cross-multiply: $4.5x = 27$, so $x = 6$ cm.}
\end{practiceQuestion}

\begin{practiceQuestion}{2-1-q14}{mc}
\begin{questionText}
What is $150\%$ of $40$?
\end{questionText}
\choiceA{$20$}
\choiceB{$40$}
\choiceC{$50$}
\choiceD{$60$}
\correctAnswer{D}
\explanation{$150\% = 1.50$. Multiply: $1.50 \times 40 = 60$. A percent greater than $100\%$ gives a result greater than the whole.}
\end{practiceQuestion}

\begin{practiceQuestion}{2-5-q22}{sa}
\begin{questionText}
Estimate a $20\%$ tip on a $\$93$ bill. Round the bill first, then calculate.
\end{questionText}
\correctAnswer{About $\$18.60$}
\explanation{Round to $\$93$. Find $10\%$: $\$9.30$. Double for $20\%$: $9.30 \times 2 = \$18.60$.}
\end{practiceQuestion}

\begin{practiceQuestion}{2-7-q28}{gmc}
\begin{questionText}
The number line below shows an estimated value and an actual value. What is the percent error?

\begin{center}
\begin{tikzpicture}[scale=0.6]
  \draw[thick,->] (-0.5,0) -- (11,0);
  \foreach \x/\lab in {0/0,2/10,4/20,6/30,8/40,10/50} {
    \draw (\x,0.15) -- (\x,-0.15) node[below,font=\small] {\lab};
  }
  % Actual at 40
  \filldraw[funGreen] (8,0) circle (4pt);
  \node[above,font=\small,funGreen!80!black] at (8,0.3) {Actual: $40$};
  % Estimated at 35
  \filldraw[funRed] (7,0) circle (4pt);
  \node[below=10pt,font=\small,funRed] at (7,-0.15) {Estimated: $35$};
  % Arrow showing difference
  \draw[thick,<->,funOrange] (7,0.7) -- (8,0.7) node[midway,above,font=\small]{$5$};
\end{tikzpicture}
\end{center}
\end{questionText}
\choiceA{$5\%$}
\choiceB{$10\%$}
\choiceC{$12.5\%$}
\choiceD{$14.3\%$}
\correctAnswer{C}
\explanation{$\frac{|35-40|}{40} \times 100 = \frac{5}{40} \times 100 = 12.5\%$.}
\end{practiceQuestion}

\begin{practiceQuestion}{2-8-q15}{mc}
\begin{questionText}
You charge $\$400$ on a credit card at $10\%$ annual interest. You pay $\$200$ right away and owe $\$200$ for $1$ year. How much total interest do you pay?
\end{questionText}
\choiceA{$\$40$}
\choiceB{$\$20$}
\choiceC{$\$10$}
\choiceD{$\$60$}
\correctAnswer{B}
\explanation{You owe $\$200$ for $1$ year. Interest $= 200 \times 0.10 = \$20$.}
\end{practiceQuestion}

\begin{practiceQuestion}{3-1-q13}{mc}
\begin{questionText}
For which value of $n$ is $|n| = 10$?
\end{questionText}
\choiceA{Only $n = 10$}
\choiceB{Only $n = -10$}
\choiceC{$n = 10$ or $n = -10$}
\choiceD{$n = 0$}
\correctAnswer{C}
\explanation{Both $|10| = 10$ and $|-10| = 10$. Two numbers have the same absolute value: the number and its opposite.}
\end{practiceQuestion}

\begin{practiceQuestion}{4-1-q08}{mc}
\begin{questionText}
Evaluate $-4y + 10$ when $y = 3$.
\end{questionText}
\choiceA{$22$}
\choiceB{$-2$}
\choiceC{$2$}
\choiceD{$-22$}
\correctAnswer{B}
\explanation{Substitute: $-4(3) + 10 = -12 + 10 = -2$.}
\end{practiceQuestion}

\begin{practiceQuestion}{4-2-q25}{sa}
\begin{questionText}
Write an expression with three terms that simplifies to $4n + 7$.
\end{questionText}
\correctAnswer{Answers vary, e.g., $n + 3n + 7$ or $5n - n + 7$}
\explanation{Any expression whose like terms combine to give $4n + 7$ is correct. For example: $n + 3n + 7 = 4n + 7$.}
\end{practiceQuestion}

\begin{practiceQuestion}{5-2-q27}{gmc}
\begin{questionText}
The table shows a pattern. What is the value of $x$ when the output is $41$?

\begin{center}
\begin{tabular}{|c|c|}
\hline
\textbf{Input ($x$)} & \textbf{Output} \\
\hline
$1$ & $8$ \\
\hline
$2$ & $11$ \\
\hline
$3$ & $14$ \\
\hline
$?$ & $41$ \\
\hline
\end{tabular}
\end{center}
\end{questionText}
\choiceA{$x = 12$}
\choiceB{$x = 11$}
\choiceC{$x = 13$}
\choiceD{$x = 14$}
\correctAnswer{A}
\explanation{The pattern is $3x + 5$. Solve $3x + 5 = 41$: subtract $5$ to get $3x = 36$, divide by $3$: $x = 12$. Check: $3(12) + 5 = 41$ \checkmark.}
\end{practiceQuestion}

\begin{practiceQuestion}{5-3-q10}{mc}
\begin{questionText}
Solve $5(n + 3) = -10$.
\end{questionText}
\choiceA{$n = -5$}
\choiceB{$n = 1$}
\choiceC{$n = -1$}
\choiceD{$n = -\dfrac{25}{5}$}
\correctAnswer{A}
\explanation{Divide by $5$: $n + 3 = -2$. Subtract $3$: $n = -5$. Check: $5(-5 + 3) = 5(-2) = -10$ \checkmark.}
\end{practiceQuestion}

\begin{practiceQuestion}{5-4-q24}{sa}
\begin{questionText}
Solve $\dfrac{5n}{6} + \dfrac{1}{3} = 3$.
\end{questionText}
\correctAnswer{$n = \dfrac{16}{5}$ or $3\dfrac{1}{5}$}
\explanation{Multiply every term by $6$: $5n + 2 = 18$. Subtract $2$: $5n = 16$. Divide by $5$: $n = \frac{16}{5} = 3\frac{1}{5}$.}
\end{practiceQuestion}

\begin{practiceQuestion}{5-5-q30}{gsa}
\begin{questionText}
The table shows ticket prices for groups. A school has $\$200$ to spend and must also pay a $\$35$ bus fee. Write and solve an inequality to find the maximum number of students $s$ who can go.

\begin{center}
\begin{tabular}{|c|c|}
\hline
\textbf{Group Size} & \textbf{Price per Student} \\
\hline
$1$--$10$ & $\$12$ \\
\hline
$11$--$25$ & $\$9$ \\
\hline
$26+$ & $\$7$ \\
\hline
\end{tabular}
\end{center}

Assume the group qualifies for the $\$9$ rate.
\end{questionText}
\correctAnswer{$9s + 35 \leq 200$; $s \leq 18$}
\explanation{Total cost: $9s + 35 \leq 200$. Subtract $35$: $9s \leq 165$. Divide by $9$: $s \leq 18.\overline{3}$. Since $s$ must be a whole number, $s \leq 18$.}
\end{practiceQuestion}

\begin{practiceQuestion}{5-6-q22}{sa}
\begin{questionText}
Solve $10 - 4x < 2$. Describe the graph.
\end{questionText}
\correctAnswer{$x > 2$; open circle at $2$, shade right}
\explanation{Subtract $10$: $-4x < -8$. Divide by $-4$ and flip: $x > 2$. Open circle at $2$, shade right.}
\end{practiceQuestion}

\begin{practiceQuestion}{6-3-q25}{sa}
\begin{questionText}
The shortest two sides of a triangle are $7$ and $9$. What is the largest whole-number length the third side can be?
\end{questionText}
\correctAnswer{$15$}
\explanation{The third side must be less than $7 + 9 = 16$. The largest whole number less than $16$ is $15$.}
\end{practiceQuestion}

\begin{practiceQuestion}{6-4-q21}{sa}
\begin{questionText}
What condition (SSS, SAS, ASA, AAS, or AAA) gives infinitely many triangles?
\end{questionText}
\correctAnswer{AAA}
\explanation{Three angles determine the shape but not the size. Infinitely many similar triangles can share the same three angle measures.}
\end{practiceQuestion}

\begin{practiceQuestion}{6-5-q11}{mc}
\begin{questionText}
Which 3D figure always has circular cross-sections when cut parallel to its base?
\end{questionText}
\choiceA{Rectangular prism}
\choiceB{Triangular pyramid}
\choiceC{Cone}
\choiceD{Cube}
\correctAnswer{C}
\explanation{A cone has a circular base, and all horizontal cuts parallel to the base produce circles of varying sizes.}
\end{practiceQuestion}

\begin{practiceQuestion}{6-6-q22}{sa}
\begin{questionText}
Two supplementary angles are in the ratio $5:4$. Find both angles.
\end{questionText}
\correctAnswer{$100°$ and $80°$}
\explanation{$5x + 4x = 180$. Combine: $9x = 180$. Solve: $x = 20$. Angles: $5(20) = 100°$ and $4(20) = 80°$.}
\end{practiceQuestion}

\begin{practiceQuestion}{7-3-q28}{gmc}
\begin{questionText}
A square has a circle inscribed inside it as shown. The side of the square is $10$ cm. What is the area of the circle? Use $\pi \approx 3.14$.

\begin{center}
\begin{tikzpicture}[scale=0.4]
  \draw[thick] (-5,-5) rectangle (5,5);
  \draw[thick,funBlue,fill=funBlue!10] (0,0) circle (5);
  \draw[thick,funRed,<->] (-5,-6) -- (5,-6);
  \node[below,funRed] at (0,-6) {$10$ cm};
  \fill (0,0) circle (2pt);
\end{tikzpicture}
\end{center}
\end{questionText}
\choiceA{$31.4$ cm$^2$}
\choiceB{$78.5$ cm$^2$}
\choiceC{$100$ cm$^2$}
\choiceD{$314$ cm$^2$}
\correctAnswer{B}
\explanation{The circle's diameter equals the side of the square: $d = 10$, so $r = 5$ cm. $A = \pi r^2 = 3.14 \times 25 = 78.5$ cm$^2$.}
\end{practiceQuestion}

\begin{practiceQuestion}{7-4-q23}{sa}
\begin{questionText}
A rectangular patio is $12$ ft by $8$ ft. A circular hot tub with a diameter of $6$ ft is placed on the patio. What is the area of the patio NOT covered by the hot tub? Use $\pi \approx 3.14$.
\end{questionText}
\correctAnswer{$67.74$ ft$^2$}
\explanation{Patio: $12 \times 8 = 96$ ft$^2$. Hot tub: $\pi r^2 = 3.14 \times 9 = 28.26$ ft$^2$. Remaining: $96 - 28.26 = 67.74$ ft$^2$.}
\end{practiceQuestion}

\begin{practiceQuestion}{7-5-q02}{mc}
\begin{questionText}
What is the surface area of a rectangular prism with length $5$ cm, width $3$ cm, and height $2$ cm?
\end{questionText}
\choiceA{$30$ cm$^2$}
\choiceB{$62$ cm$^2$}
\choiceC{$46$ cm$^2$}
\choiceD{$31$ cm$^2$}
\correctAnswer{B}
\explanation{$SA = 2(5 \times 3) + 2(5 \times 2) + 2(3 \times 2) = 30 + 20 + 12 = 62$ cm$^2$.}
\end{practiceQuestion}

\begin{practiceQuestion}{7-6-q03}{mc}
\begin{questionText}
What is the volume of a cube with side length $5$ m?
\end{questionText}
\choiceA{$15$ m$^3$}
\choiceB{$25$ m$^3$}
\choiceC{$125$ m$^3$}
\choiceD{$150$ m$^3$}
\correctAnswer{C}
\explanation{$V = s^3 = 5^3 = 125$ m$^3$.}
\end{practiceQuestion}

\begin{practiceQuestion}{8-1-q21}{sa}
\begin{questionText}
A random sample of $50$ apples from an orchard found that $6$ were bruised. If the orchard has $3{,}000$ apples, about how many are bruised?
\end{questionText}
\correctAnswer{$360$}
\explanation{$\frac{6}{50} = 12\%$. Then $12\%$ of $3{,}000 = 360$.}
\end{practiceQuestion}

\begin{practiceQuestion}{8-3-q22}{sa}
\begin{questionText}
Team A: mean $60$, MAD $4$. Team B: mean $70$, MAD $3$. Is the difference in means meaningful? Explain briefly.
\end{questionText}
\correctAnswer{Yes; the means are about $2.5$ to $3.3$ MADs apart}
\explanation{Difference $= 10$. Using MAD of $4$: $\frac{10}{4} = 2.5$. Using MAD of $3$: $\frac{10}{3} \approx 3.3$. Both are greater than $2$, so the difference is meaningful.}
\end{practiceQuestion}

\begin{practiceQuestion}{8-4-q05}{mc}
\begin{questionText}
Data set: $20, 22, 25, 27, 30, 32, 35$. What is the IQR?
\end{questionText}
\choiceA{$15$}
\choiceB{$10$}
\choiceC{$7$}
\choiceD{$5$}
\correctAnswer{B}
\explanation{$Q_1 = 22$, $Q_3 = 32$. IQR $= 32 - 22 = 10$.}
\end{practiceQuestion}

\begin{practiceQuestion}{8-6-q22}{sa}
\begin{questionText}
A survey of $60$ people: $18$ chose option A, $24$ chose option B, $18$ chose option C. Find the angle for option B in a circle graph.
\end{questionText}
\correctAnswer{$144°$}
\explanation{$\frac{24}{60} = 0.40 = 40\%$. Angle $= 0.40 \times 360° = 144°$.}
\end{practiceQuestion}

\begin{practiceQuestion}{9-4-q11}{mc}
\begin{questionText}
A bag contains $5$ yellow and $5$ purple marbles. Lisa picks a marble, records the color, and replaces it $40$ times. She gets yellow $24$ times. Which statement is correct?
\end{questionText}
\choiceA{The model predicts $P(\text{yellow}) = 0.5$, and the experiment gives $P(\text{yellow}) = 0.6$, so there is a small difference}
\choiceB{The model is wrong because $24 \neq 20$}
\choiceC{The model predicts $P(\text{yellow}) = 0.6$}
\choiceD{The experiment proves yellow is more likely than purple}
\correctAnswer{A}
\explanation{Predicted: $\frac{5}{10} = 0.5$. Experimental: $\frac{24}{40} = 0.6$. The small difference is expected from natural variation.}
\end{practiceQuestion}

\begin{practiceQuestion}{9-6-q20}{sa}
\begin{questionText}
A coin is flipped and a die is rolled. What is the probability of getting tails and a $1$? Write your answer as a fraction.
\end{questionText}
\correctAnswer{$\frac{1}{12}$}
\explanation{$P(\text{tails}) = \frac{1}{2}$, $P(1) = \frac{1}{6}$. $P = \frac{1}{2} \times \frac{1}{6} = \frac{1}{12}$.}
\end{practiceQuestion}

\begin{practiceQuestion}{9-7-q12}{mc}
\begin{questionText}
How does increasing the number of trials in a simulation affect the results?
\end{questionText}
\choiceA{The results become less reliable}
\choiceB{The experimental probability gets closer to the theoretical probability}
\choiceC{The results stay exactly the same}
\choiceD{The simulation takes less time}
\correctAnswer{B}
\explanation{More trials produce more reliable estimates. The experimental probability approaches the theoretical value as trials increase.}
\end{practiceQuestion}
